\documentclass[11pt,a4paper]{article}
\usepackage{fontspec,amsmath,amssymb,amsthm,microtype}
\setmainfont{lmroman10-regular.otf}[BoldFont=lmroman10-bold.otf,ItalicFont=lmroman10-italic.otf,BoldItalicFont=lmroman10-bolditalic.otf]
\usepackage[margin=25mm]{geometry}
\usepackage[unicode,hidelinks]{hyperref}
\hypersetup{pdfauthor={LangevinLean Project}}
\setlength{\emergencystretch}{3em}
\setlength{\parskip}{0.4em}
\setlength{\parindent}{0pt}
\allowdisplaybreaks[2]
\raggedbottom
\title{\texorpdfstring{The local Euler defect}{The local Euler defect}}
\author{LangevinLean Project}
\date{}
\begin{document}\maketitle
Draft mathematical exposition; these notes may contain mistakes.


Let \(H\) be a real Banach space, \(g:H\to H\) a Lipschitz map, and \(Y,B\) curves satisfying the stated drift integral equation. All time integrals are Bochner integrals. Define \(D_h=\int_0^h(g(Y_t)-g(Y_0))dt\). Integral linearity gives \(Y_h=Y_0-hg(Y_0)+B_h-D_h\). The norm bound follows by integrating the Lipschitz estimate for \(g(Y_t)-g(Y_0)\).



For the application \(H=L^2(P;\mathbb R^d)\), the norms are mean-square norms. The assumed drift and noise bounds give \(\|Y_t-Y_0\|\le t\sqrt{Ld}+\sqrt{2d}\sqrt t\). Their integrals have orders \(h^2\) and \(h^{3/2}\). Under \(Lh\le2\), the resulting defect is at most \((33/20)Lh\sqrt{hd}\). Stationarity and the score bound remain hypotheses of this calculation and are constructed separately.


\section{\texorpdfstring{Definition of the local defect}{Definition of the local defect}}


For \(g:H\to H\), \(Y:\mathbb R\to H\), and \(h\in\mathbb R\), define the drift omitted when Euler freezes \(g\) at the initial point:

\[D_h(g,Y):=\int_0^h\bigl(g(Y_t)-g(Y_0)\bigr)\,dt.\]

Integrate the difference between the evolving drift and its frozen value. This defines the defect itself; the estimates below will be consequences of this definition.



The integrand vanishes at time zero, so this defect measures drift variation rather than the drift magnitude itself. In a space of random variables, the integral is a Bochner integral in that space and its norm is an \(L^2\) error. The definition integrates only drift variation; the additive input occurs separately in the integral equation.


\section{\texorpdfstring{Separate the exact and frozen drift}{Separate the exact and frozen drift}}


If \(t\mapsto g(Y_t)\) is integrable on the interval between \(0\) and \(h\), then

\[D_h(g,Y)=\int_0^h g(Y_t)\,dt-h\,g(Y_0).\]

Expand the definition. Both the time-dependent drift and the constant function \(g(Y_0)\) are integrable, so linearity separates their difference into two integrals.



The integral of the constant vector is \((h-0)g(Y_0)=h\,g(Y_0)\), yielding the identity.



Subtracting the constant inside the integral is legitimate because a constant vector is integrable on a finite interval. For negative \(h\), the oriented integral changes sign consistently on both terms. Later error estimates assume \(h\ge0\) so that time-dependent upper bounds can be integrated in their usual order.


\section{\texorpdfstring{An exact step is Euler minus its defect}{An exact step is Euler minus its defect}}


Assume the preceding drift integrability and the endpoint integral equation

\[Y_h-Y_0=-\int_0^h g(Y_t)\,dt+B_h.\]

Then the exact reference step has the form

\[Y_h=Y_0-h\,g(Y_0)+B_h-D_h(g,Y).\]

Solve the assumed integral equation for \(Y_h\), then replace the defect by the difference between the exact drift integral and its frozen approximation.



Cancel the two frozen-drift terms. The resulting equality is precisely the Euler update with the signed remainder \(-D_h\).



Substitute \(\int_0^hg(Y_t)dt=hg(Y_0)+D_h\) into the endpoint equation. The minus sign in front of the original drift integral gives \(-hg(Y_0)-D_h\). This sign convention is retained when comparing the numerical and exact updates, where subtracting the reference changes \(-D_h\) into \(+D_h\).


\section{\texorpdfstring{Continuity supplies drift integrability}{Continuity supplies drift integrability}}


Let \(H\) be a normed additive group, \(g:H\to H\) be \(L\)-Lipschitz with \(L\ge0\), and \(Y:\mathbb R\to H\) be continuous. For every real \(h\), the function \(t\mapsto g(Y_t)\) is interval integrable from \(0\) to \(h\).



A Lipschitz map is continuous, so \(g\circ Y\) is continuous. A continuous function is integrable on every compact interval, proving the claim.



The image of a compact time interval under a continuous path is compact. A continuous drift is bounded on that image, so the composition has finite integral norm there. In the lifted random-vector setting the same argument uses continuity in \(L^2\), not a pathwise boundedness claim about every sample trajectory.


\section{\texorpdfstring{The square-root integral}{The square-root integral}}


For every \(h\ge0\),

\[\int_0^h\sqrt t\,dt=\frac23h\sqrt h.\]

Write \(\sqrt t=t^{1/2}\) and use the power-integral formula. The exponent satisfies \(1/2>-1\), so the formula applies at the zero endpoint.



The contribution at zero vanishes. Since \(h\ge0\), the other endpoint is \(h^{1/2+1}=h^{1/2}h=\sqrt h\,h\).



Divide by \(1/2+1=3/2\) and rearrange the factors to obtain \((2/3)h\sqrt h\).



The integrand has an integrable singularity in its derivative at zero, but the function \(\sqrt t\) itself is continuous there. The primitive \((2/3)t^{3/2}\) vanishes at zero. Consequently the Brownian contribution to an integrated Lipschitz increment is exactly proportional to \(h^{3/2}\), with the constant \(2/3\) retained for the final numerical estimate.


\section{\texorpdfstring{Bound a reference increment}{Bound a reference increment}}


Let \(h\ge0\) and suppose that, for every \(0\le s\le h\),

\[\|g(Y_s)\|\le G,\qquad \|B_s\|\le N\sqrt s,\qquad Y_s-Y_0=-\int_0^s g(Y_r)\,dr+B_s.\]

Then, for each \(0\le t\le h\),

\[\|Y_t-Y_0\|\le tG+N\sqrt t.\]

First bound the drift integral on \([0,t]\) by its constant norm bound:

\[\left\|\int_0^t g(Y_s)\,ds\right\|\le G\,|t-0|=Gt,\]

where \(t\ge0\) removes the absolute value.



To justify the constant bound under that integral, note that \(s\in(0,t]\) implies \(0\le s\le t\le h\). Hence the assumed bound on \([0,h]\) applies to every such \(s\).



Insert the integral equation and apply the triangle inequality:

\[\|Y_t-Y_0\|\le\left\|\int_0^t g(Y_s)\,ds\right\|+\|B_t\|.\]

Negation leaves the norm of the drift integral unchanged.



Apply the two bounds to obtain \(Gt+N\sqrt t\), and commute the factors \(G\) and \(t\).



For each \(t\), the drift integral runs only over \([0,t]\), which lies in the assumed control interval \([0,h]\). Its norm is at most \(\int_0^tGds=tG\). Adding the noise estimate gives the stated bound; no independence between drift and noise is used, because this is a norm inequality before taking any probabilistic expectation.


\section{\texorpdfstring{Integrate the Lipschitz increment bound}{Integrate the Lipschitz increment bound}}


Assume \(h\ge0\), \(g\) is \(L\)-Lipschitz with \(L\ge0\), and for \(0\le t\le h\) the same integral equation and bounds \(\|g(Y_t)\|\le G\), \(\|B_t\|\le N\sqrt t\) hold. Then

\[\|D_h(g,Y)\|\le L\left(\frac{h^2G}{2}+\frac23h\sqrt h\,N\right).\]

For \(0<t\le h\), Lipschitz continuity gives

\[\|g(Y_t)-g(Y_0)\|\le L\|Y_t-Y_0\|.\]

Multiply the increment estimate by \(L\ge0\) to obtain the pointwise majorant

\[\|g(Y_t)-g(Y_0)\|\le L(tG+N\sqrt t).\]

The scalar majorant \(t\mapsto L(tG+N\sqrt t)\) is continuous, hence integrable on \([0,h]\).



Expand \(D_h\). The norm estimate for a Bochner integral, together with the integrable majorant, gives

\[\|D_h\|\le\int_0^h L(tG+N\sqrt t)\,dt.\]

Pull out constants and split the two integrable summands. Evaluate

\[\int_0^h t\,dt=\frac{h^2}{2},\qquad \int_0^h\sqrt t\,dt=\frac23h\sqrt h.\]

This leaves the two drift and noise contributions in the claimed expression.



Rearrange the scalar factors to place the drift term first and the noise term second.



The two terms have distinct origins and scales: \(Lh^2G/2\) comes from the deterministic accumulation of drift, and \((2/3)LNh^{3/2}\) comes from the square-root growth of the additive input. Substitution of the stationary moment bounds can therefore be made separately in these two terms.


\section{\texorpdfstring{The exact rational constant}{The exact rational constant}}


The numerical factor in the corrected Dalalyan–Karagulyan estimate satisfies

\[\frac76\sqrt2\le\frac{33}{20}=1.65.\]

Use \((\sqrt2)^2=2\) and \(\sqrt2\ge0\). Comparing squares proves the bound, since \(49/18\le1089/400\).



Both sides are nonnegative, so their squares may be compared without changing the order. After squaring, the difference is the rational positive number \(1089/400-49/18\). This proves the rational upper bound on \(\sqrt2\).


\section{\texorpdfstring{The local error with stationary moment scales}{The local error with stationary moment scales}}


Let \(h,d\in\mathbb R\) with \(h,d\ge0\), let \(g\) be \(L\)-Lipschitz with \(L\ge0\), and suppose \(Lh\le2\). Assume that, for every \(0\le t\le h\),

\[Y_t-Y_0=-\int_0^t g(Y_s)\,ds+B_t,\quad \|g(Y_t)\|\le\sqrt{Ld},\quad \|B_t\|\le\sqrt{2d}\sqrt t.\]

Then

\[\boxed{\|D_h(g,Y)\|\le\frac{33}{20}Lh\sqrt{hd}.}\]

In \(L^2\), the drift bound corresponds to a score second moment at most \(Ld\); here that bound remains a hypothesis.



Compare the nonnegative quantities \(h\sqrt{Ld}\) and \(\sqrt2\sqrt{hd}\) by squaring. It suffices to show \(h^2Ld\le2hd\).



Replace each squared square root by its nonnegative argument. Multiply \(Lh\le2\) by \(hd\ge0\), obtaining \(h^2Ld\le2hd\), and thus

\[h\sqrt{Ld}\le\sqrt2\sqrt{hd}.\]

Split products under the square roots and rearrange:

\[h\sqrt h\sqrt{2d}=h\sqrt2\sqrt{hd}.\]

The nonnegativity of \(h\) and \(2\) justifies the square-root product identities.



Apply the two-term local-defect bound with \(G=\sqrt{Ld}\) and \(N=\sqrt{2d}\):

\[\|D_h\|\le L\left(\frac{h^2\sqrt{Ld}}2+\frac23h\sqrt h\sqrt{2d}\right).\]

Bound the first term using \(h\sqrt{Ld}\le\sqrt2\sqrt{hd}\), and rewrite the second with the preceding identity. Since \(1/2+2/3=7/6\),

\[\|D_h\|\le L\frac76\sqrt2\,h\sqrt{hd}.\]

Multiply \((7/6)\sqrt2\le33/20\) by the nonnegative factor \(Lh\sqrt{hd}\). This gives the claimed constant without numerical rounding.



The condition \(Lh\le2\) converts the drift contribution, initially of order \(h^2\sqrt{Ld}\), into the same \(Lh\sqrt{hd}\) scale as the Brownian contribution. The sum of coefficients is \(\sqrt2(1/2+2/3)=(7/6)\sqrt2\). Applying the proved upper bound \(33/20\) is the only numerical enlargement in this local-defect estimate.


\end{document}
